\documentclass[11pt]{article}
\usepackage[T1]{fontenc}
\usepackage{amsmath,amssymb,amsthm,mathtools}
\usepackage[margin=1in]{geometry}
\usepackage{hyperref}
\newtheorem{theorem}{Theorem}
\newtheorem{lemma}[theorem]{Lemma}
\newtheorem{proposition}[theorem]{Proposition}
\newcommand{\R}{\mathbb R}
\newcommand{\C}{\mathbb C}
\newcommand{\PF}{\operatorname{PF}}
\newcommand{\sech}{\operatorname{sech}}
\renewcommand{\theequation}{GC\arabic{equation}}
\title{An exact correction of the printed Guinand--Weil archimedean comparison}
\author{}
\date{19 September 2026}
\begin{document}
\maketitle

This note concerns the equation labelled
\texttt{eq:guinand-weil-arch-map} and its expanded version
\texttt{eq:guinand-weil-residual-complete} in the programme's
\emph{Original Six Readers}.  The local equations, cosine convention and
finite-part subtraction are retained literally.  The result is a direct
calculation of their exact difference, with no use of RH or of a zero sum.
The incoming note \emph{Exact correction to the archived Guinand--Weil
comparison}, 19 September 2026, identified this difference.  Here every
interchange and endpoint is proved, the test-domain assertion is corrected,
and the entire exponential family is evaluated.  The result corrects the
local transcription/comparison; it makes no assertion against Guinand's
published theorem \cite{guinand1948} or Weil's published formula
\cite{weil1952}.

\section{The original functions and the exact domain inclusion}

For $c>0$, define $\mathcal D_c$ to consist of complex continuous functions
$f:[0,\infty)\to\C$ with a finite partition
$0=a_0<a_1<\cdots<a_m<\infty$ such that $f$ is $C^2$ on each open piece,
$f''$ is absolutely integrable on every piece, and $f'$ has finite
one-sided values at the finite endpoints.  On the final unbounded piece
require
\begin{equation}
 |f(x)|+|f'(x)|+|f''(x)|\le C_f e^{-cx}
 \quad (x\ge X_f).
 \label{gc:decay}
\end{equation}
At a derivative jump use its two one-sided values, and use their average
as its assigned representative.  This representative does not change an
integral.  Put
\begin{equation}
 \mathcal D_0^{\exp}=\bigcup_{c>0}\mathcal D_c,
 \qquad \mathcal D_W^{\exp}=\bigcup_{c>1/2}\mathcal D_c.
 \label{gc:domains}
\end{equation}
These are vector spaces: for addition take the common refinement of the
two finite partitions and the smaller of the two decay exponents.

The original definitions are
\begin{align}
 g_f(t)&=\sqrt{\frac2\pi}\int_0^\infty f(x)\cos(tx)\,dx,
 &&t\ge0,\label{gc:cosine}\\
 2H_f&=\int_0^\infty f(x)
       \left(\frac1x-\frac{e^{-x/2}}{\sinh x}\right)dx,
 \label{gc:H}\\
 K(x)&=\frac{e^{|x|/2}}{|e^x-e^{-x}|},\qquad x\ne0,
 \label{gc:K}\\
 P_f&=\lim_{\Lambda\to\infty}
 \left\{\int_\R(1-e^{-\Lambda|x|})f(|x|)K(x)\,dx
              -f(0)\log\Lambda\right\},
 \label{gc:PF}\\
 \mathscr L_f&=\sqrt{\frac2\pi}\int_0^\infty
             g_f(t)\log\frac{t}{2\pi}\,dt,
 \label{gc:L}\\
 W_f&=2\int_0^\infty f(x)e^{-x/2}\,dx-P_f-f(0)\log(2\pi).
 \label{gc:W}
\end{align}
In \eqref{gc:PF}, $|x|f(|x|)K(x)\to f(0)/2$, so the original
Weil subtraction $2\beta(0)\log\Lambda$ is exactly
$f(0)\log\Lambda$; this factor is not adjusted.
The archived equality was $W_f=2H_f+\mathscr L_f$.

\begin{proposition}[Exact inclusion in the recorded common domain]
For every $f\in\mathcal D_W^{\exp}$ and $1<p\le2$, the function $f$
satisfies the Guinand test conditions recorded in the source:
it is an indefinite integral of its locally integrable derivative,
$f(x)\to0$, and $xf'(x)\in L^p(0,\infty)$.
The map
\begin{equation}
 E:\mathcal D_W^{\exp}\longrightarrow
       \mathcal W_{\rm even}^{G,p},\qquad
 E(f)(x)=\frac{f(|x|)}{\sqrt{2\pi}},
 \qquad f(t)=\sqrt{2\pi}\,E(f)(t)\ (t\ge0)
 \label{gc:E}
\end{equation}
is an injective linear map into the recorded Guinand--Weil common test
space, and the displayed restriction is its inverse on its image.
Moreover
\begin{equation}
 \int_\R E(f)(x)e^{itx}\,dx=g_f(t),\qquad
 T^{1/2}f(T)\longrightarrow0.
 \label{gc:transform-map}
\end{equation}
\end{proposition}
\begin{proof}
The integrability of $f''$ and finite derivative traces make $f'$ bounded
on each bounded piece.  Continuity of $f$ cancels the boundary values in
the fundamental theorem on adjacent pieces.  Thus
$f(x)=f(0)+\int_0^x f'(u)du$ on every bounded interval.
The exponential bound gives the limit at infinity and the $L^p$ assertion;
near zero and each finite break the derivative is bounded.
The even extension is continuous and piecewise continuously
differentiable.  Its derivative at $0$ has one-sided values
$-f'(0+)/\sqrt{2\pi}$ and $f'(0+)/\sqrt{2\pi}$; at $\pm a_j$ it
has precisely the corresponding reflected derivative jumps.  Thus it
satisfies the recorded Weil condition (A), including the averaged
representatives.  Choose a decay exponent $c>1/2$ in \eqref{gc:decay} and
put $b=c-1/2>0$.  Then both $E(f)$ and its piecewise derivative are
$O(e^{-(1/2+b)|x|})$, which is condition (B).
The two expressions in \eqref{gc:E} compose to the identity.  Splitting
the integral at zero gives
$2/\sqrt{2\pi}=\sqrt{2/\pi}$ times the cosine integral, proving
\eqref{gc:transform-map}.  Its last limit follows directly from
\eqref{gc:decay}.  These are the common-domain conditions recorded from
Guinand's Theorems 3 and $3'$ and Weil's (A), (B)
\cite[pp.~112--115]{guinand1948}\cite[Collected Papers II, pp.~52--58]{weil1952}.
\end{proof}

The same expression $E(f)$ is a defined even Fourier test for
$f\in\mathcal D_0^{\exp}$ and satisfies the same transform identity,
but the inclusion into Weil's (B) requires the faster exponent.
For example $f(x)=e^{-x/4}$ belongs to $\mathcal D_0^{\exp}$, whereas
\begin{equation}
 e^{(1/2+b)x}|E(f)(x)|
   =\frac{e^{(1/4+b)x}}{\sqrt{2\pi}}\longrightarrow\infty
   \quad (b>0).
 \label{gc:domain-counterexample}
\end{equation}
Consequently the incoming note's phrase ``exponential decay'' alone does
not prove its claimed inclusion in the original common domain.  Equations
\eqref{gc:domains}--\eqref{gc:transform-map} give the exact inclusion and
inverse.  All functional identities below hold on
$\mathcal D_0^{\exp}$, and their application as a correction within the
recorded common domain is restricted to $\mathcal D_W^{\exp}$.

\section{Cosine decay, inversion at zero, and the logarithmic functional}

\begin{lemma}[Jump-retaining decay and the endpoint inversion]
For $f\in\mathcal D_0^{\exp}$ set
$\Delta_j=f'(a_j+)-f'(a_j-)$ and
$B_f=|f'(0+)|+\sum_{j=1}^m|\Delta_j|+\int_0^\infty|f''(x)|dx$,
where $f''$ is the piecewise derivative.  For $t>0$,
\begin{equation}
 g_f(t)=-\frac{\sqrt{2/\pi}}{t^2}
 \left(f'(0+)+\sum_{j=1}^m\Delta_j\cos(ta_j)
                   +\int_0^\infty f''(x)\cos(tx)dx\right).
 \label{gc:IBP}
\end{equation}
In particular $g_f$ and $g_f(t)\log(t/(2\pi))$ are absolutely integrable,
and
\begin{equation}
 \sqrt{\frac2\pi}\int_0^\infty g_f(t)dt=f(0).
 \label{gc:inversion}
\end{equation}
\end{lemma}
\begin{proof}
First integrate $f(x)\cos(tx)$ by parts on all pieces up to a finite
endpoint $R$ in the last piece.  Continuity makes the internal $f$
boundary terms cancel.  The term at zero contains $\sin0=0$, and the
term at $R$ tends to zero.  This gives
$\int f\cos=-t^{-1}\int f'\sin$.
On the second integration, $\sin(tx)=-t^{-1}(\cos(tx))'$.
The endpoint at zero contributes $f'(0+)/t$; the two values at $a_j$
give $\Delta_j\cos(ta_j)/t$; the endpoint at infinity vanishes.
Multiplication by the first factor $-1/t$ proves \eqref{gc:IBP}.
Thus
$|g_f(t)|\le\sqrt{2/\pi}\min(\|f\|_1,B_f/t^2)$.
This proves integrability of the transform and of its logarithmic weight
both at zero and infinity.

For $r>0$, absolute convergence permits the exact Laplace--cosine
calculation
\begin{equation}
 \sqrt{\frac2\pi}\int_0^\infty e^{-rt}g_f(t)dt
  =\frac2\pi\int_0^\infty f(x)\frac{r}{r^2+x^2}dx
  =\frac2\pi\int_0^\infty\frac{f(ru)}{1+u^2}du.
 \label{gc:Poisson}
\end{equation}
Indeed $\int_0^\infty e^{-rt}\cos(tx)dt$
is the real part of $(r-ix)^{-1}$ and equals $r/(r^2+x^2)$;
the double integral is bounded by $\|f\|_1/r$.
The function $f$ is bounded.  Dominated convergence in the right-hand
integral as $r\downarrow0$ gives $f(0)$, since
$\int_0^\infty(1+u^2)^{-1}du=\pi/2$.
Dominated convergence by $|g_f|$ in the left-hand integral proves
\eqref{gc:inversion} without an unproved endpoint inversion assertion.
\end{proof}

\begin{theorem}[The exact cosine-log functional]
For $f\in\mathcal D_0^{\exp}$,
\begin{equation}
 \boxed{\displaystyle
 \mathscr L_f=-\int_0^\infty\frac{f(x)-f(0)e^{-x}}{x}dx
                      -f(0)\log(2\pi).}
 \label{gc:log-formula}
\end{equation}
\end{theorem}
\begin{proof}
Set $h(x)=f(x)-f(0)e^{-x}$.  It belongs to $\mathcal D_0^{\exp}$,
$h(0)=0$, and $h(x)=O(x)$ near zero because its right derivative is
bounded there.  Its exponential decay gives
\begin{equation}
 \int_0^\infty\frac{|h(x)|}{x}dx<\infty,
 \qquad \int_0^\infty g_h(t)dt=0.
 \label{gc:h}
\end{equation}
For $t>0$ the Frullani identity used here has an elementary proof:
\begin{equation}
 \int_0^\infty\frac{e^{-r}-e^{-tr}}{r}dr
 =\int_1^t\left(\int_0^\infty e^{-ur}dr\right)du
 =\int_1^t\frac{du}{u}=\log t.
 \label{gc:Frullani}
\end{equation}
The interval $[1,t]$ is oriented if $t<1$; absolute integrability on the
bounded $u$ interval justifies the exchange.  For fixed $t$ the integrand
in $r$ has constant sign.  Therefore every truncated integral
$\ell_{\epsilon,R}(t)=\int_\epsilon^R(e^{-r}-e^{-tr})dr/r$
satisfies $|\ell_{\epsilon,R}(t)|\le|\log t|$.

At $0<\epsilon<R<\infty$, Fubini and \eqref{gc:h} give
\begin{align}
 \sqrt{\frac2\pi}\int_0^\infty g_h(t)\ell_{\epsilon,R}(t)dt
 &=-\sqrt{\frac2\pi}\int_\epsilon^R\frac1r
                        \int_0^\infty e^{-rt}g_h(t)dt\,dr\notag\\
 &=-\frac2\pi\int_\epsilon^R\int_0^\infty
                           \frac{h(x)}{r^2+x^2}dx\,dr.
 \label{gc:truncated-log}
\end{align}
For the first exchange, $2\|g_h\|_1\log(R/\epsilon)$ is a finite
majorant.  For inserting the cosine transform in the second line use
$\int_\epsilon^R r^{-1}\int_0^\infty e^{-rt}dt\,dr\,\|h\|_1<\infty$.
For the untruncated double integral, Tonelli applied to the absolute value
gives exactly
\begin{equation}
 \int_0^\infty\int_0^\infty\frac{|h(x)|}{r^2+x^2}dx\,dr
   =\frac\pi2\int_0^\infty\frac{|h(x)|}{x}dx<\infty.
 \label{gc:absolute-double}
\end{equation}
The left side of \eqref{gc:truncated-log} tends to the logarithmic
integral by the majorant $|g_h(t)\log t|$ from the lemma.  The right side
tends by \eqref{gc:absolute-double} to $-\int h(x)/x\,dx$.

For the retained reference $e^{-x}$ its transform is
$\sqrt{2/\pi}/(1+t^2)$.  The substitution $t=1/u$ in the absolutely
convergent integral yields
$\int_0^\infty\log t\,(1+t^2)^{-1}dt=0$.
Hence its $\mathscr L$ value is $-\log(2\pi)$.
Adding $f(0)$ times this reference to the result for $h$ proves
\eqref{gc:log-formula}, including its constant.
\end{proof}

\section{The original finite part and the defect map}

\begin{theorem}[Evaluation with the original regulator]
For $f\in\mathcal D_0^{\exp}$ the limit in \eqref{gc:PF} exists and
\begin{equation}
 \boxed{\displaystyle
 P_f=\int_0^\infty\left(
       \frac{f(x)e^{x/2}}{\sinh x}-\frac{f(0)e^{-x}}{x}\right)dx.}
 \label{gc:PF-formula}
\end{equation}
The displayed combined integrand is absolutely integrable.
\end{theorem}
\begin{proof}
Write $A_f(x)=f(x)e^{x/2}/\sinh x-f(0)e^{-x}/x$ for $x>0$.
The expansions
$e^{x/2}/\sinh x=x^{-1}+1/2+O(x)$ and
$e^{-x}/x=x^{-1}-1+O(x)$, together with
$f(x)=f(0)+f'(0+)x+o(x)$, show
$A_f(x)\to f'(0+)+3f(0)/2$ at zero.
At infinity its two terms decay exponentially.  Thus $A_f\in L^1$.
Evenness changes the full-line regulated integral into
$\int_0^\infty(1-e^{-\Lambda x})f(x)e^{x/2}/\sinh x\,dx$,
with no residual factor of two.  Subtract and add the reference to obtain
the exact finite-$\Lambda$ identity
\begin{align}
 &\int_\R(1-e^{-\Lambda|x|})f(|x|)K(x)dx-f(0)\log\Lambda\notag\\
 &\qquad=\int_0^\infty(1-e^{-\Lambda x})A_f(x)dx
              +f(0)\log(1+\Lambda^{-1}).
 \label{gc:finite-regulator}
\end{align}
Here the reference integral is
$\int_0^\infty(e^{-x}-e^{-(1+\Lambda)x})dx/x=\log(1+\Lambda)$
by \eqref{gc:Frullani}.  Dominated convergence by $|A_f|$, and
$\log(1+\Lambda^{-1})\to0$, prove \eqref{gc:PF-formula}.
\end{proof}

\begin{theorem}[Exact comparison and receiving map]
Define
\begin{equation}
 d(x)=2e^{-x/2}-\sech(x/2)
     =\frac{2e^{-3x/2}}{1+e^{-x}},\qquad
 D:\mathcal D_0^{\exp}\to\C,\quad D(f)=\int_0^\infty f(x)d(x)dx.
 \label{gc:D}
\end{equation}
Then the original expressions satisfy the exact identity
\begin{equation}
 \boxed{W_f-(2H_f+\mathscr L_f)=D(f).}
 \label{gc:defect}
\end{equation}
In particular the map
\begin{equation}
 T:\mathcal D_0^{\exp}\to\C^3,\quad T(f)=(W_f,H_f,\mathscr L_f),
 \qquad (1,-2,-1)\,T=D
 \label{gc:typed}
\end{equation}
is the exact typed relation between the two presentations.  Its restriction
along the inclusion $\mathcal D_W^{\exp}\hookrightarrow\mathcal D_0^{\exp}$
is the corrected common-domain comparison.
The defect is strictly positive on every nonzero real nonnegative test.
\end{theorem}
\begin{proof}
The kernel in \eqref{gc:H} tends to $1/2$ at zero and is $O(1/x)$ at
infinity, proving absolute convergence of $H_f$.  All other terms were
proved convergent above.  Substitute \eqref{gc:log-formula} and
\eqref{gc:PF-formula}; to avoid separating divergent terms, first combine
their integrands on $[\eta,R]$ and then pass to the limits.  The
$f(0)\log(2\pi)$ constants and the reference subtractions cancel exactly,
as do the $f(x)/x$ terms.  The remaining kernel is
\begin{equation}
 2e^{-x/2}+\frac{e^{-x/2}-e^{x/2}}{\sinh x}
 =2e^{-x/2}-\frac1{\cosh(x/2)}=d(x).
 \label{gc:kernel-algebra}
\end{equation}
It is bounded at zero and bounded by $2e^{-3x/2}$ at infinity.
Thus the limit is $D(f)$, proving \eqref{gc:defect} and \eqref{gc:typed}.
The second expression for $d$ in \eqref{gc:D} is strictly positive.
A continuous nonnegative nonzero $f$ is positive on an interval, so its
integral against $d$ is strictly positive.
\end{proof}

For an explicit receiving functional retaining all old terms, set
\begin{align}
 H_f^{\rm rec}&=H_f+\tfrac12D(f),\label{gc:receiving-H}\\
 2H_f^{\rm rec}
 &=\int_0^\infty f(x)
  \left(\frac1x-\frac{e^{x/2}}{\sinh x}+2e^{-x/2}\right)dx,
 \qquad W_f=2H_f^{\rm rec}+\mathscr L_f.
 \label{gc:receiving-kernel}
\end{align}
The kernel in this integral tends to $3/2$ at zero.  Equations
\eqref{gc:receiving-H}--\eqref{gc:receiving-kernel} follow by the already
proved identity and \eqref{gc:kernel-algebra}; they explicitly record the
map from the archived $H_f$.  They do not rename the archived $H_f$ or
identify $H_f^{\rm rec}$ with an uninspected historical transcription.

\section{The full exponential family and the original counterexample}

For $a>0$ let $f_a(x)=e^{-ax}$.  Its cosine transform, obtained by taking
the real part of $(a-it)^{-1}$, is
\begin{equation}
 g_{f_a}(t)=\sqrt{\frac2\pi}\frac{a}{a^2+t^2}.
 \label{gc:exponential-transform}
\end{equation}
In particular $f_a\in\mathcal D_0^{\exp}$ for all $a>0$ and
$f_a\in\mathcal D_W^{\exp}$ for $a>1/2$.
Put $z=a/2+1/4$ and $w=a/2+3/4$.  We obtain
\begin{align}
 P_{f_a}&=-\log2-\psi(z),\label{gc:exp-PF}\\
 2H_{f_a}&=\psi(w)+\log(2/a),\label{gc:exp-H}\\
 \mathscr L_{f_a}&=\log(a/(2\pi)),\label{gc:exp-L}\\
 W_{f_a}&=\frac{2}{a+1/2}+\psi(z)-\log\pi,
 \label{gc:exp-W}\\
 2H_{f_a}+\mathscr L_{f_a}&=\psi(w)-\log\pi,
 \label{gc:exp-right}\\
 D(f_a)&=\frac{2}{a+1/2}+\psi(z)-\psi(w)
       =\psi(z+1)-\psi(w)>0.
 \label{gc:exp-D}
\end{align}
Here $\psi=\Gamma'/\Gamma$ has its usual unscaled argument.

To verify every constant in these identities, use the primary special
function formula \cite[5.9.12]{dlmf}
\begin{equation}
 \psi(v)=\int_0^\infty
 \left(\frac{e^{-t}}t-\frac{e^{-vt}}{1-e^{-t}}\right)dt,
 \qquad v>0.
 \label{gc:digamma}
\end{equation}
Its integrand is $O(1)$ at zero and exponentially integrable at infinity.
In \eqref{gc:PF-formula}, substitute $t=2x$ to get the combined
integrand $e^{-zt}/(1-e^{-t})-e^{-t/2}/t$.
Adding and subtracting $e^{-t}/t$ inside the combined integral gives
$-\psi(z)$ and
$\int_0^\infty(e^{-t}-e^{-t/2})dt/t=-\log2$ by
\eqref{gc:Frullani}.  This proves \eqref{gc:exp-PF}.
The same substitution in \eqref{gc:H} gives
$e^{-at/2}/t-e^{-wt}/(1-e^{-t})$; the difference from
\eqref{gc:digamma} integrates to $\log(2/a)$, proving
\eqref{gc:exp-H}.
In the transform integral put $t=au$.  The constant term integrates to
$\log(a/(2\pi))$, while the $\log u$ term is zero by $u\mapsto1/u$.
This proves \eqref{gc:exp-L}.  The remaining identities follow by
substitution, using the exact recurrence
$\psi(z+1)=\psi(z)+1/z$ \cite[5.5.2]{dlmf}; here
$1/z=2/(a+1/2)$.  Positivity follows independently from \eqref{gc:D}.

For the incoming test $a=3/2$ one has $z=1$, $w=3/2$.
The special values
$\psi(1)=-\gamma_E$ and $\psi(1/2)=-\gamma_E-2\log2$,
and the recurrence at $1/2$, give
\cite[5.4.12--13 and 5.5.2]{dlmf}
\begin{align}
 P_{f_{3/2}}&=\gamma_E-\log2,&
 2H_{f_{3/2}}&=2-\gamma_E-\log3,\notag\\
 \mathscr L_{f_{3/2}}&=\log\frac3{4\pi},&
 W_{f_{3/2}}&=1-\gamma_E-\log\pi,\notag\\
 2H_{f_{3/2}}+\mathscr L_{f_{3/2}}
   &=2-\gamma_E-\log\pi-2\log2.
 \label{gc:counterexample-values}
\end{align}
Consequently
\begin{equation}
 \boxed{D(f_{3/2})=2\log2-1>0.}
 \label{gc:counterexample}
\end{equation}
There is also a direct evaluation independent of digamma values:
\begin{equation}
 D(f_{3/2})=2\int_0^\infty\frac{e^{-3x}}{1+e^{-x}}dx
 =2\int_0^1\frac{u^2}{1+u}du
 =2\left(\frac12-1+\log2\right).
 \label{gc:direct-example}
\end{equation}
Strict positivity also follows from
$\log2=\int_1^2dt/t>1/2$, without a numerical approximation.
Thus the contradiction to the archived equality occurs inside its actual
recorded common domain, and does not depend on the repaired inclusion for
more slowly decaying functions.

\section{Exact supersession and provenance}

Equation \eqref{gc:defect} supersedes the claimed zero-defect equality at
the following two archived labels:
\texttt{eq:guinand-weil-arch-map} and
\texttt{eq:guinand-weil-residual-complete}.
The proof that the latter follows from the two displayed local scalar
formulas cannot validate both scalar transcriptions with the printed
definitions: subtracting them gives the contradicted equality.  The
calculation here locates and evaluates the discrepancy exactly, while
preserving both displayed functionals and all their original factors.
It does not determine from this discrepancy alone which published formula
was miscopied.  That would require comparing the historical formulas
themselves, rather than using the local assertion that they were literal.

The source snapshot is
\texttt{06\_Foundations\_and\_Original\_Six\_Readers.tex}, SHA--256
\begin{center}\small\ttfamily
e653a19b9bfc7b4788550e59de42367d52ed666598e46540d040e06d7362cad0.
\end{center}
Its relevant original lines are 201784--201795 (Guinand domain and
transform), 201798--201843 (extension), 201953--202030 (expanded false
comparison), 202639--202660 (Weil conditions and regulator), and
202842--202906 (printed comparison).  The incoming correction is
\texttt{06\_PRINTED\_ARCHIMEDEAN\_CORRECTION.md}, SHA--256
\begin{center}\small\ttfamily
6d0f28606fe98d8ab49f25b54d94377e1481ce83a511ce8189bfebfba629dcee.
\end{center}
The source files are preserved unchanged.  The three functional identities
and the admissible exponential example are established here with complete
proofs; no zero-location conclusion, mixed-support control, or arithmetic
allocation follows merely from this correction.

\begin{thebibliography}{9}
\bibitem{guinand1948}
A.~P.~Guinand, ``A Summation Formula in the Theory of Prime Numbers,''
\emph{Proceedings of the London Mathematical Society}, second series
\textbf{50} (1948), 107--119.
\href{https://doi.org/10.1112/plms/s2-50.2.107}{doi:10.1112/plms/s2-50.2.107}.
The historic attribution and theorem numbers are retained from the original
reader; the functional correction is proved directly above, rather than
assuming that reader's scalar transcriptions.
\bibitem{weil1952}
A.~Weil, ``Sur les formules explicites de la th\'eorie des nombres premiers,''
\emph{Communications du S\'eminaire Math\'ematique de l'Universit\'e de Lund},
supplement dedicated to Marcel Riesz (1952), 252--265; reprinted in
\emph{Collected Papers}, vol.~II, Springer, 1979, 48--61.
Conditions (A), (B) and the finite-part convention used here are the ones
displayed in the source snapshot.
\bibitem{dlmf}
F.~W.~J. Olver, A.~B. Olde Daalhuis, D.~W. Lozier, B.~I. Schneider,
R.~F. Boisvert, C.~W. Clark, B.~R. Miller, B.~V. Saunders, H.~S. Cohl,
and M.~A. McClain, eds., \emph{NIST Digital Library of Mathematical
Functions}, Chapter 5.
\href{https://dlmf.nist.gov/5.9.E12}{Equation 5.9.12},
\href{https://dlmf.nist.gov/5.5.E2}{equation 5.5.2}, and
\href{https://dlmf.nist.gov/5.4}{equations 5.4.12--13};
accessed 19 September 2026.
\end{thebibliography}
\end{document}
